\section{Evaluating the weights with bounded rational computation}
\label{sec:evaluation}

The block integrals in Section~\ref{sec:ports} are definitions of weights,
not oracles supplied with the input. We now construct their evaluators.
The central step expresses a close approximation to an exponential
integral as the volume of a rational polytope. We check its separation
oracle, encoding lengths, and inner and outer geometric guarantees.
These checks also explain why a free-coordinate interval of length
$R+1$ causes only polynomial dependence on $\log(R+1)$.

\subsection{Magnitude bounds before numerical approximation}

\begin{lemma}[Integral and point-weight bounds]\label{lem:weight-bounds}
For every block and every port argument in its box,
\begin{equation}\label{eq:integral-bounds}
 e^{-N^{50000}}\le\widetilde I_G(u)\le e^K.
\end{equation}
For the weight $h$ on the full paired slice of
Proposition~\ref{prop:port-model},
\begin{equation}\label{eq:weight-log-bound}
 |\log h(a,b)|\le N^{60000}.
\end{equation}
All rational coefficients needed to define or evaluate its factors
have polynomial encoding length in the input and in the query.
\end{lemma}

\begin{proof}
The integral's upper bound is its free-box volume,
$(R+1)^{|\mathcal F_G|}\le e^K$, since smoothing uses probability kernels.
For the lower bound put $u^c_{iD}=c(D)$ and $z^c_e=c_e$ on free edges.
The tree solve gives $X_1=X_2=c$, so the prefix-slack argument of
Lemma~\ref{lem:port-penalties} gives
$V_G(u^c,z^c)\le2N^4S$. For each free coordinate choose
$z^*_e=\max\{c_e,1/4\}$. It lies in $[0,R+1]$ with at least $1/4$
clearance from its boundary and differs from $c_e$ by at most $1/4$.
On the cube of sup radius $1/8$ about $z^*$, and for every smoothing
displacement, the port and free-coordinate displacements from
$(u^c,z^c)$ are at most $N^{21}W+b_0+1$. Therefore
\begin{equation}\label{eq:near-center-penalty}
 \lambda V_G(u-v,z)
 \le2\lambda N^4S+L_0(N^{21}W+b_0+1)<N^{5000}.
\end{equation}
Integrating over this cube and over the whole probability kernel gives
$\widetilde I_G(u)\ge4^{-|\mathcal F_G|}e^{-N^{5000}}
\ge e^{-N^{50000}}$. This proof includes $\mathcal F_G=\varnothing$.

For the rank factors, choose any common integral base $x^*$, which
exists by the feasibility and integrality results of Section~\ref{sec:polyhedra}.
Its projected totals $v^*_i$ satisfy the rank constraint and are within
$N^{10}W$ of the center totals by~\eqref{eq:part-total-range}.
For any nonnegative vector $v$ on a rank index set, its rank penalty is
\[
 v(E_r)+r(E_r)-2m_r(v)
  =\max_{A\subseteq E_r}
       \{v(A)-v(E_r\setminus A)+r(E_r)-2r(A)\}.
\]
Every affine function on the right has coordinate slopes $\pm1$.
Thus this penalty is 1-Lipschitz in the $1$-norm. It vanishes at $v_i^*$,
so throughout a port box it is at most $N^{23}W$, independently of the
numerical value of $R$. Multiplication by $H$ still gives a fixed power
of $N$. For each component, $t_C$ vanishes at the center and has
absolute value at most $N^{23}W$ throughout the boxes. Consequently
all Gaussian and absolute-value exponents in~\eqref{eq:block-factor}
have absolute value bounded by fixed powers below $N^{40000}$; the
auxiliary quadratic is smaller than one. Summing these bounds and
\eqref{eq:integral-bounds} over at most $n$ blocks proves
\eqref{eq:weight-log-bound}.

The center, endpoints, scales, and tree coefficients have polynomial
encoding length. Ranks of unions and contracted ranks require ordinary
rank queries. For the minima defining $m_r$ and the maxima defining
$V_G$, rational modular terms are handled by clearing their common
denominator and applying Proposition~\ref{prop:oracle-tools}. The binary
length of that denominator is at most the sum of the individual
denominator lengths. Thus the size parameter is polynomial in bits;
no bound on its numerical value is needed.
\end{proof}

\subsection{A rational volume problem}

We use the following form of the classical randomized volume theorem
\cite{DFK1991,DyerFrieze1991}. In dimension $d$, a full-dimensional
rational polytope supplied by strong separation, a polynomial bound
on each inequality's encoding length, and rational inner and outer
ball guarantees admits a relative volume approximation with success
probability at least $3/4$. Its oracle and bit complexity is polynomial
in $d$, these encoding lengths, and the reciprocal error. The radii
enter through their binary lengths. We use the finite-precision volume
algorithm in the rational oracle model of
\cite[Section~2]{DyerFrieze1991}. Its polynomial random-bit bound is
recorded in \cite[Section~6]{DyerFrieze1991}; finite-grid sampling within
cubes and its error analysis appear in
\cite[algorithm Step~2 and its finite-grid error analysis]{DFK1991}.
Finite uniform integer choices can be generated from unbiased bits by
rejection using fewer than two expected trials. The resulting expected
polynomial bound includes bit operations, rational arithmetic, and
oracle calls; we cap and amplify this procedure below. Strong separation
supplies the weak membership required in the volume theorem.

We first round each volume input by a rational change of coordinates.
Apply Lemma~\ref{lem:center} to the full-dimensional
polytope and send its returned simplex vertices to $0,e_1,\ldots,e_d$.
The image contains the standard simplex and lies in $[-2,2]^d$,
because its nonconstant barycentric coordinates have absolute value
at most two. The standard simplex contains a Euclidean ball about
its barycenter of radius at least $(d+1)^{-2}$. Translating that
barycenter to zero and scaling by $(d+1)^2$ therefore give an inner
unit ball and an outer ball of polynomial radius. The affine map,
its determinant, and the inverse images of rational separation queries
have polynomial encoding length. Correcting the returned volume by
the determinant is exact rational arithmetic. The required optimization
uses the rational separation theorem \cite{GLS1988}, through
Proposition~\ref{prop:oracle-tools}. This supplies a rounded input
even when the original geometric aspect ratio is numerically large.

Fix a block and a rational port vector $u$ in its box. Let $m$ be the
number of its port arguments, including fixed length-zero ports, and
$f=|\mathcal F_G|$ its number of free variables. Length-zero ports are
omitted only from the count slice; they remain arguments of the smoothing
convolution.
Replace its $6m$ uniform variables by $q_{ja}\in[0,1]$, putting
\[
 v_j(q)=\frac{b_0}{3}\sum_{a=1}^6(q_{ja}-1/2),\qquad
 \mathcal B=[0,R+1]^f\times[0,1]^{6m},\qquad
 D(z,q)=\lambda V_G(u-v(q),z).
\]
The normalized uniform density has been absorbed by this unit-interval
parametrization, so there is no further Jacobian factor:
\begin{equation}\label{eq:exponential-box-integral}
 \widetilde I_G(u)=\int_{\mathcal B}e^{-D(x)}\dd x,
 \qquad x=(z,q).
\end{equation}
Here $D$ is nonnegative, convex, and piecewise affine. Its value and an
attaining affine piece can be computed exactly at rational points:
submodular minimization returns a maximizing subset for each of the
two maxima in $V_G$, and their sum is such a piece.

For $0<\rho<1$, define
\begin{equation}\label{eq:volume-lift-scales}
 E_0=N^{51000}+\lceil\log_2(100/\rho)\rceil,
 \qquad d_0=\lceil1000E_0^2/\rho\rceil.
\end{equation}
The ceiling of the logarithm can be found by integer comparisons of
powers of two, without numerical logarithms. Consider the polytope
\begin{equation}\label{eq:volume-lift}
 \mathcal Q=\{(x,y):x\in\mathcal B,\ y\in\R_{\ge0}^{d_0},\quad
                   D(x)+d_0y_a\le d_0\ (1\le a\le d_0)\}.
\end{equation}
Each inequality involving $D$ means all the affine-piece inequalities
in its finite maximum representation. This makes $\mathcal Q$ a
convex rational polytope, although its inequalities are not enumerated.
Integrating over its $y$-fiber gives the exact identity
\begin{equation}\label{eq:fiber-volume}
 \vol(\mathcal Q)=\int_{\mathcal B}(1-D(x)/d_0)_+^{d_0}\dd x.
\end{equation}

\begin{lemma}[Accuracy and guarantees for the volume lift]
\label{lem:volume-lift}
The polytope~\eqref{eq:volume-lift} has dimension $f+6m+d_0$, a
polynomial facet-encoding bound, and an exact polynomial-time strong
separation oracle. It contains a known rational ball of radius $1/32$
and has an outer ball of polynomial encoding length. Furthermore,
\begin{equation}\label{eq:volume-lift-accuracy}
 (1-\rho/10)\widetilde I_G(u)
       \le\vol(\mathcal Q)\le\widetilde I_G(u).
\end{equation}
All polynomial bounds are in the original input length, the encoding
lengths of $u$ and $\rho$, $N$, and $\rho^{-1}$, with absolute fixed exponents.
\end{lemma}

\begin{proof}
The pointwise upper inequality follows from $1-t\le e^{-t}$. For
$0\le D\le E_0$, expansion of $-\log(1-t)$, or integration of its
derivative, gives
\[
 0\le-d_0\log(1-D/d_0)-D
 \le\frac{D^2}{2d_0(1-D/d_0)}<\rho/100.
\]
The part of~\eqref{eq:exponential-box-integral} where $D>E_0$ has
integral at most $e^{K-E_0}$. Dividing by the lower bound in
\eqref{eq:integral-bounds} makes its relative contribution at most
\[
 \exp(K+N^{50000}-E_0)<\rho/100.
\]
These two errors prove~\eqref{eq:volume-lift-accuracy}.

For separation, first test the box inequalities and $y\ge0$. Compute
$D(x)$ and choose an index $a$ with largest $y_a$. If
$D(x)+d_0y_a>d_0$, the attaining affine piece of $D$ gives a rational
hyperplane separating the point from $\mathcal Q$. Otherwise all
inequalities hold. The coefficient bounds follow from the integral
tree map, the rank-answer bit bound, the rational scales, and the
input port encodings. At a rational query, clearing denominators for
the two submodular minimizations takes polynomial bit work as in
Lemma~\ref{lem:weight-bounds}. Each inequality, rather than the number
of possible subsets, is the encoding parameter.

For the inner ball choose
\[
 x^*=(z^*,q^*),\qquad z^*_e=\max\{c_e,1/4\},\qquad
 q^*_{ja}=1/2,\qquad y^*_a=1/4.
\]
All box coordinates have clearance at least $1/4$. A Euclidean
perturbation of radius $1/32$ changes each free coordinate by at most
$1/32$ and each $v_j$ by at most $2b_0/32$. By
\eqref{eq:port-lipschitz} its change in $D$ is at most
$L_0\max\{1,2b_0\}/32<1$. The near-center estimate
\eqref{eq:near-center-penalty} therefore gives $D<N^{5000}+1<d_0/2$
throughout this ball. Also $0<y_a\le9/32$, so
$D+d_0y_a<25d_0/32<d_0$. This proves the inner guarantee, including
$f=0$. Every point of $\mathcal Q$ has $x\in\mathcal B$ and
$0\le y_a\le1$. If $d=f+6m+d_0$, the ball centered at the origin
with rational radius $d(R+2)$ is therefore an outer ball. Its binary
encoding length is polynomial.

Finally, the numerical value of $d_0$ is polynomial in $N$ and
$\rho^{-1}$: the exponents in~\eqref{eq:volume-lift-scales} are fixed.
Thus even its explicitly added coordinates can be written and
processed within one fixed polynomial bound. The original capacity
appears only in the box endpoint $R+1$ and the encoding lengths just
verified.
\end{proof}

\subsection{Positive rational outputs and deterministic caps}

\begin{proposition}[A bounded weight evaluator]\label{prop:weight-evaluator}
Given a profile in the entire slice $\mathcal S$, rational
$0<\rho,\beta<1$, and a rational power $s\in[0,1]$ with denominator
$q_s$, there is an evaluator using fresh random bits and returning a strictly positive rational
$\widehat h_s$ such that
\[
 \PP\bigl[(1-\rho)h^s\le\widehat h_s\le(1+\rho)h^s\bigr]
       \ge1-\beta.
\]
Its output length, number of rank queries, and other bit operations
are bounded on every execution by fixed polynomials in $N$, the
input encoding lengths, $\rho^{-1}$, $\log(\beta^{-1})$, and $q_s$.
In particular an equally spaced schedule with a polynomial number of
phases has polynomially bounded denominators. Moreover $J$ admits
a deterministic positive rational approximation $\widehat J$ with
$|\log(\widehat J/J)|\le N^{-7}$ in polynomial work.
\end{proposition}

\begin{proof}
First evaluate a single smoothed block integral. Round $\mathcal Q$
by the rational affine construction above, apply the volume theorem
with relative error $\rho/10$, and restore the exact determinant factor.
Use the same parameter $\rho$ in~\eqref{eq:volume-lift-scales}. By
Lemma~\ref{lem:volume-lift}, its successful output is a relative
$\rho/2$ approximation to $\widetilde I_G$.

For completeness, all resource bounds here are bounds on actual
finite-bit executions. Let $B$ be a furnished polynomial upper bound
on the expected combined number of bit operations and oracle calls
of one volume run, including the cost of its separation queries.
If the chosen version already has a deterministic bound, use that
bound instead. Interrupt the run after $16B$ operations; Markov's
inequality increases the failure probability by at most $1/16$.
The resulting success probability is at least $11/16$. On interruption,
a declared failure (including a zero denominator), or an invalid or
nonpositive output, substitute a fixed positive rational. Repeat independently
$O(1+\log\beta^{-1})$ times and take the median. The binomial tail
bound gives failure probability at most $\beta$. Cap construction,
query lengths, arithmetic, and random-bit generation are all included.
Clip each output to
\[
 [2^{-2N^{50000}},\ 2^{2N^{50000}}].
\]
This interval has polynomial-bit endpoints and strictly contains all
successful relative approximations from~\eqref{eq:integral-bounds}.
Clipping therefore retains the success guarantee and ensures bounded
positive outputs even on failed or interrupted runs.

Evaluate all rank penalties and the other rational exponents exactly,
and allocate relative error $\rho/(100(n+1))$ and failure allowance
$\beta/(n+1)$ to each of the at most $n$ block integrals. Evaluate the
remaining exponentials and all rational arithmetic to another
$\rho/(100(n+1))$ relative error per factor. The log bounds in
Lemma~\ref{lem:weight-bounds} ensure that these relative accuracies
require only polynomial absolute precision. One elementary
implementation uses the exponential Taylor series and factorial tail
bounds: an exponent of magnitude at most $D_*$ and absolute precision
$2^{-k}$ require polynomially many terms and bits in $D_*+k+1$.
Negative exponents can be evaluated by inversion of a positive
exponential approximation. Here $D_*$ is a fixed polynomial in $N$,
not in $R$.

Multiplying the factors gives a positive relative approximation to
$h$. Before taking any power, clip this product to
$[2^{-3N^{60000}},2^{3N^{60000}}]$. This interval contains $h$ by
\eqref{eq:weight-log-bound}, so clipping cannot increase its relative
error; it also gives valid bisection bounds even after failed block
evaluations. On the simultaneous success event the sum of the factor
logarithmic errors is below $\rho/10$, after reducing the indicated
constants if necessary. To obtain $h^s$ for $s=a_s/q_s$, apply bisection to the
$q_s$th root and integer powering to $a_s$, with positive rational
endpoints whose binary exponents are bounded by
$3N^{60000}$. Each comparison is an integer-power comparison, and
$a_s\le q_s$; the number of bisections, lengths of operands, and
bit work are polynomial in $q_s$, $N$, and $\log\rho^{-1}$.
Taking a power in $[0,1]$ cannot increase the preceding logarithmic
error. Approximate the root to logarithmic error at most
$\rho/[100(q_s+1)]$ before taking the $a_s$th power; the resulting
additional logarithmic error is at most $\rho/100$. This proves the
required relative bound. For $s=0$ return one exactly.
All steps have prescribed finite bounds; clipping the final output
to $[2^{-3N^{60000}},2^{3N^{60000}}]$ gives the same positivity and
encoding guarantee on every failure event. This proves the asserted
evaluator without conditioning any future computation on its success.

To approximate $J$, let $k$ be the number of nonroot components, so
\[
 J=\left(\sum_{j\in\Z}e^{-j^2/M^2}\right)^k,
 \qquad 0\le k\le n.
\]
For $k=0$ return one. Otherwise truncate at the integer
$B_J=MN^{10}$, a polynomial in $N$. The omitted two-sided tail is at
most
\[
 2\int_{B_J}^{\infty}e^{-x^2/M^2}\dd x
 \le\frac{M^2}{B_J}e^{-B_J^2/M^2}
 =MN^{-10}e^{-N^{20}}.
\]
The sum is at least one. Approximate each of its $2B_J+1$ retained
terms to absolute error at most
$N^{-7}/[100(n+1)(2B_J+1)]$, keeping the term at zero exactly one
and the remaining terms nonnegative. The elementary exponential
procedure above is bounded since each exponent has magnitude at most
$N^{20}$. Adding the terms and raising to the integer power $k$
gives a positive rational result whose logarithmic error is below
$N^{-7}$. No random choices or numerical grid of length $R$ occurs.
\end{proof}
